\documentclass[10pt,a4paper]{amsart}
\usepackage{amsmath,amssymb,amsthm,graphicx}
\usepackage{verbatim}
\newenvironment{reference}{\comment}{\endcomment}
%\newenvironment{reference}{}{}
\newenvironment{slogan}{\comment}{\endcomment}
\newenvironment{history}{\comment}{\endcomment}
% For commutative diagrams we use Xy-pic
\usepackage[all]{xy}

% We use 2cell for 2-commutative diagrams.
\xyoption{2cell}
% Theorem environments.
%
\theoremstyle{plain}
\newtheorem{theorem}[subsection]{Theorem}
\newtheorem{proposition}[subsection]{Proposition}
\newtheorem{lemma}[subsection]{Lemma}

\theoremstyle{definition}
\newtheorem{definition}[subsection]{Definition}
\newtheorem{example}[subsection]{Example}
\newtheorem{exercise}[subsection]{Exercise}
\newtheorem{situation}[subsection]{Situation}

\theoremstyle{remark}
\newtheorem{remark}[subsection]{Remark}
\newtheorem{remarks}[subsection]{Remarks}

\numberwithin{equation}{subsection}

% Macros
%
\def\lim{\mathop{\mathrm{lim}}\nolimits}
\def\colim{\mathop{\mathrm{colim}}\nolimits}
\def\Spec{\mathop{\mathrm{Spec}}}
\def\Hom{\mathop{\mathrm{Hom}}\nolimits}
\def\Ext{\mathop{\mathrm{Ext}}\nolimits}
\def\SheafHom{\mathop{\mathcal{H}\!\mathit{om}}\nolimits}
\def\SheafExt{\mathop{\mathcal{E}\!\mathit{xt}}\nolimits}
\def\Sch{\mathit{Sch}}
\def\Mor{\mathop{\mathrm{Mor}}\nolimits}
\def\Ob{\mathop{\mathrm{Ob}}\nolimits}
\def\Sh{\mathop{\mathit{Sh}}\nolimits}
\def\NL{\mathop{N\!L}\nolimits}
\def\CH{\mathop{\mathrm{CH}}\nolimits}
\def\proetale{{pro\text{-}\acute{e}tale}}
\def\etale{{\acute{e}tale}}
\def\QCoh{\mathit{QCoh}}
\def\Ker{\mathop{\mathrm{Ker}}}
\def\Im{\mathop{\mathrm{Im}}}
\def\Coker{\mathop{\mathrm{Coker}}}
\def\Coim{\mathop{\mathrm{Coim}}}

% Boxtimes
%
\DeclareMathSymbol{\boxtimes}{\mathbin}{AMSa}{"02}

%
% Macros for moduli stacks/spaces
%
\def\QCohstack{\mathcal{QC}\!\mathit{oh}}
\def\Cohstack{\mathcal{C}\!\mathit{oh}}
\def\Spacesstack{\mathcal{S}\!\mathit{paces}}
\def\Quotfunctor{\mathrm{Quot}}
\def\Hilbfunctor{\mathrm{Hilb}}
\def\Curvesstack{\mathcal{C}\!\mathit{urves}}
\def\Polarizedstack{\mathcal{P}\!\mathit{olarized}}
\def\Complexesstack{\mathcal{C}\!\mathit{omplexes}}
% \Pic is the operator that assigns to X its picard group, usage \Pic(X)
% \Picardstack_{X/B} denotes the Picard stack of X over B
% \Picardfunctor_{X/B} denotes the Picard functor of X over B
\def\Pic{\mathop{\mathrm{Pic}}\nolimits}
\def\Picardstack{\mathcal{P}\!\mathit{ic}}
\def\Picardfunctor{\mathrm{Pic}}
\def\Deformationcategory{\mathcal{D}\!\mathit{ef}}
\usepackage{cleveref}
\crefname{originallemma}{original lemma}{original lemmas}
\Crefname{originallemma}{Original lemma}{Original lemmas}
\makeatletter
\newlabel{lemma-finite-type-points-morphism}{{lemma-finite-type-points-morphism}{source}{}{sourcepointer-1}{}}
\newlabel{lemma-finite-type-points-morphism@cref}{{[originallemma][0][]lemma-finite-type-points-morphism}{[1][0][]source}{}{}{}}
\makeatother
\crefname{originallemma}{original lemma}{original lemmas}
\Crefname{originallemma}{Original lemma}{Original lemmas}
\makeatletter
\newlabel{lemma-identify-finite-type-points}{{lemma-identify-finite-type-points}{source}{}{sourcepointer-2}{}}
\newlabel{lemma-identify-finite-type-points@cref}{{[originallemma][0][]lemma-identify-finite-type-points}{[1][0][]source}{}{}{}}
\makeatother
\crefname{originallemma}{original lemma}{original lemmas}
\Crefname{originallemma}{Original lemma}{Original lemmas}
\makeatletter
\newlabel{lemma-immersion-locally-finite-type}{{lemma-immersion-locally-finite-type}{source}{}{sourcepointer-3}{}}
\newlabel{lemma-immersion-locally-finite-type@cref}{{[originallemma][0][]lemma-immersion-locally-finite-type}{[1][0][]source}{}{}{}}
\makeatother
\crefname{originalreference}{original reference}{original references}
\Crefname{originalreference}{Original reference}{Original references}
\makeatletter
\newlabel{stacks-properties-remark-stack-structure-locally-closed-subset}{{stacks-properties-remark-stack-structure-locally-closed-subset}{source}{}{sourcepointer-4}{}}
\newlabel{stacks-properties-remark-stack-structure-locally-closed-subset@cref}{{[originalreference][0][]stacks-properties-remark-stack-structure-locally-closed-subset}{[1][0][]source}{}{}{}}
\makeatother
\title{Every nonempty locally closed subset of the space of a stack meets the set of points of finite type}
\author{The Stacks Project Authors}
\date{stacks-project:stacks-morphisms.tex:local-snapshot-20261012; GNU FDL 1.2 or later}
\begin{document}
\maketitle
\noindent\textit{Editorial note.} All delivered mathematical text is the author's own, copied from the pinned Stacks Project chapter source sources/stacks-algebraic-stacks-morphisms/pinned-original-source.tex; no statement or proof step was rewritten, reordered or removed.\par\medskip
\noindent\textit{Editorial note.} Preamble provenance: the Stacks Project keeps one single shared preamble for every chapter; this card takes it directly from the staged sources/stacks-algebraic-stacks-morphisms/preamble.tex as three verbatim spans (lines 8-12, 14-18 and 154-217), so no derived file was written and nothing was added or redefined. The wrapper class and the base packages are supplied by the assembler.\par\medskip
\noindent\textit{Editorial note.} Bibliography provenance: neither the counted statement nor its proof contains a citation (this lane rescans the whole span before writing the card), so the card carries no bibliography and no entry was reconstructed.\par\medskip
\noindent\textit{Editorial note.} Selection provenance: this card was opened under the file-based control contract (CONTROL.md), with the mechanical pre-card gate run on its statement and proof interval before the card was written.\par\medskip
\noindent\textit{Allowed standard prerequisites (editorial).}\begin{itemize}
\item The construction of the reduced induced substack structure on a locally closed subset. Invoked at the start of the counted proof (source lines 3780-3783) as a cross-chapter pointer. Turns the given subset into an algebraic stack whose space is the subset.
\item The statement that an immersion is locally of finite type. Invoked inside the counted proof (source lines 3784-3785). Allows points of finite type of the substack to be compared with those of the ambient stack.
\item The comparison of points of finite type under a morphism. Invoked inside the counted proof (source lines 3786-3788). Places the points of finite type of the substack among those of the stack inside the subset.
\item The description of points of finite type through closed points of affine presentations. Invoked at the end of the counted proof (source lines 3789-3793). Produces a point of finite type on the substack from a closed point of an affine presentation.
\end{itemize}\medskip
\ifdefined\hypertarget\hypertarget{sourcepointer-1}{}\fi
\noindent\textit{Source reference.} Original lemma lemma-finite-type-points-morphism: sources/stacks-algebraic-stacks-morphisms/pinned-original-source.tex, source lines 3728--3728. Navigation pointer only: the counted text refers to the Stacks item lemma-finite-type-points-morphism, which is an existing labeled result of this same chapter and is not part of this card. The pointer carries the target's own label as its number and its definition line as its span; its content is not reproduced here.\par\medskip
\ifdefined\hypertarget\hypertarget{sourcepointer-2}{}\fi
\noindent\textit{Source reference.} Original lemma lemma-identify-finite-type-points: sources/stacks-algebraic-stacks-morphisms/pinned-original-source.tex, source lines 3711--3711. Navigation pointer only: the counted text refers to the Stacks item lemma-identify-finite-type-points, which is an existing labeled result of this same chapter and is not part of this card. The pointer carries the target's own label as its number and its definition line as its span; its content is not reproduced here.\par\medskip
\ifdefined\hypertarget\hypertarget{sourcepointer-3}{}\fi
\noindent\textit{Source reference.} Original lemma lemma-immersion-locally-finite-type: sources/stacks-algebraic-stacks-morphisms/pinned-original-source.tex, source lines 3514--3514. Navigation pointer only: the counted text refers to the Stacks item lemma-immersion-locally-finite-type, which is an existing labeled result of this same chapter and is not part of this card. The pointer carries the target's own label as its number and its definition line as its span; its content is not reproduced here.\par\medskip
\ifdefined\hypertarget\hypertarget{sourcepointer-4}{}\fi
\noindent\textit{Source reference.} Original reference stacks-properties-remark-stack-structure-locally-closed-subset: sources/stacks-algebraic-stacks-morphisms/pinned-original-source.tex, source lines 3783--3783. Navigation pointer only: the counted text refers to the Stacks item stacks-properties-remark-stack-structure-locally-closed-subset, which belongs to another chapter of the project and is not part of this card. The pointer carries the target's own label as its number because this chapter's staged snapshot does not contain the other chapter's tag table; its content is not reproduced here.\par\medskip
\begin{lemma}
\label{lemma-enough-finite-type-points}
Let $\mathcal{X}$ be an algebraic stack.
For any locally closed subset $T \subset |\mathcal{X}|$ we have
$$
T \not = \emptyset
\Rightarrow
T \cap \mathcal{X}_{\text{ft-pts}} \not = \emptyset.
$$
In particular, for any closed subset $T \subset |\mathcal{X}|$ we
see that $T \cap \mathcal{X}_{\text{ft-pts}}$ is dense in $T$.
\end{lemma}

\begin{proof}
Let $i : \mathcal{Z} \to \mathcal{X}$ be the reduced induced substack
structure on $T$, see
Properties of Stacks,
Remark \ref{stacks-properties-remark-stack-structure-locally-closed-subset}.
An immersion is locally of finite type, see
Lemma \ref{lemma-immersion-locally-finite-type}.
Hence by
Lemma \ref{lemma-finite-type-points-morphism}
we see
$\mathcal{Z}_{\text{ft-pts}} \subset \mathcal{X}_{\text{ft-pts}} \cap T$.
Finally, any nonempty affine scheme $U$ with a smooth morphism towards
$\mathcal{Z}$ has at least one closed point, hence $\mathcal{Z}$ has at least
one finite type point by
Lemma \ref{lemma-identify-finite-type-points}.
The lemma follows.
\end{proof}

\end{document}
